% Three-policy-lane view of P-CONSENT. The three policies are the primary
% structure: pi^DR (day-ahead), pi^ch (every dt), pi^neg (per arrival), each
% on its own clock. Vertical arrows are the within-day dependency (FSL ->
% charging feasible set -> negotiation options); the dashed arrow is the sole
% cross-day link (surviving buffer). Not a timeline: see Sec. 4.
% Requires \usetikzlibrary{positioning, fit, backgrounds, arrows.meta}.
\begin{figure*}[t]
\centering
\newcommand{\ltitle}{\bfseries}%
\newcommand{\leqn}{\scriptsize\itshape\color{black!55}}%
\begin{tikzpicture}[
  font=\footnotesize,
  every node/.style={transform shape},
  lane/.style={rectangle, rounded corners=3pt, draw=black!55, thick,
               align=left, inner sep=7pt, text width=0.60\textwidth,
               minimum height=1.55cm},
  clock/.style={rectangle, rounded corners=3pt, draw=black!40, thick,
                fill=white, align=center, inner sep=4pt, text width=1.9cm,
                minimum height=1.55cm, font=\scriptsize},
  dep/.style={-{Latex[length=2.2mm]}, very thick, black!75},
  back/.style={-{Latex[length=2.2mm]}, very thick, dashed, black!75},
  deplab/.style={font=\scriptsize\itshape, text=black!60, align=center, fill=white, inner sep=1pt}
]
% ============ LANE 1: pi^DR (day-ahead) ============
\node[clock, fill=blue!8] (cDR) {\ltitle Day-ahead\\(end of $d{-}1$)\\[2pt]\leqn once/day};
\node[lane, fill=blue!5, right=0.35cm of cDR] (DR) {%
{\ltitle Layer $\pidr$: demand-response commitment.}\ \ Forecast the surviving
buffer, then commit a defensible peak cap $\FSL_d$ and earn a capacity payment.
Fixed before any day-$d$ arrival.\\[2pt]
{\leqn $\psihat\!\to\!\buffhat{d}\!\to\!\Omeps$; set
$\FSL_d(\lambda)$;\ earn $\Cinc{d}$}};

% ============ LANE 2: pi^ch (every dt) ============
\node[clock, fill=teal!9, below=1.7cm of cDR] (cCH) {\ltitle Every $\dt$\\[2pt]\leqn continuous};
\node[lane, fill=teal!6, right=0.35cm of cCH] (CH) {%
{\ltitle Layer $\pich$: real-time charging.}\ \ Set charge/discharge rates at
least cost under the committed FSL; bank only into spare capacity; during an
event (if called) discharge to hold load under $\FSL_d$.\\[2pt]
{\leqn minimize s.t.\--;
bank if $\Ptot{t}\!<\!\min(\Pmaxhat,\text{cap})$; slack $\vio{t}$ via $\Cpen{d}$}};

% ============ LANE 3: pi^neg (per arrival) ============
\node[clock, fill=orange!10, below=1.7cm of cCH] (cNEG) {\ltitle Per arrival\\/ departure\\[2pt]\leqn event-driven};
\node[lane, fill=orange!6, right=0.35cm of cNEG] (NEG) {%
{\ltitle Layer $\pineg$: negotiation \& persistence.}\ \ Offer each driver
compensated flexibility and a banking price $\pbank{v}{k}$; at departure settle
energy, banking, and realized event/peak help; carry the surviving buffer over.\\[2pt]
{\leqn price by $\Uflex$$\Rightarrow\!\Eneg,\tdep$; settle
$\ucost{k}{v}$, $\Ubuf,\Udc$--;
carry $\survived{v}{k+1}$}};

% ============ WITHIN-DAY DEPENDENCY (top -> down) ============
\draw[dep] (DR.south) -- (CH.north)
  node[deplab, midway] {$\FSL_d$ constrains charging feasible set};
\draw[dep, <->] (CH.south) -- (NEG.north)
  node[deplab, midway] {discharge capability sets offerable flexibility};

% ============ CROSS-DAY LINK (bottom -> up, dashed) ============
\draw[back] (NEG.east) -- ++(0.7,0)
  node[midway, right, font=\scriptsize\itshape, text=black!70, align=center]
    {surviving buffer\\ $\buff{d}\!\to\!\buffhat{d+1}$}
  |- (DR.east);

  

\end{tikzpicture}
\caption{The three policies of P-CONSENT, each acting on its own clock:
$\pidr$ commits the firm service level a day ahead, $\pich$ sets charging and
discharging rates every $\dt$, and $\pineg$ negotiates flexibility and settles
payments at each arrival and departure. Solid vertical arrows are the sole
within-day dependency, the committed FSL constrains the charging feasible set,
which in turn sets the flexibility the negotiation can offer; this is a
dependency, not a temporal order, since the charging and negotiation epochs
interleave continuously through the day (the problem definition (Section~\ref{sec:formulation})). The dashed
arrow is the only coupling between days: the buffer that survives a driver's
trip becomes tomorrow's deliverable buffer and sharpens the forecast that sets
the next commitment.}
\label{fig:processflow}
\end{figure*}
